We focus or analysis on the United States over the period from 1982 until 2025 and adjust the timeframe analyzed based on the desired scenarios.\footnote{Sencarios are the Great Financial Crisis in 2007-2009 and Covid-19 in 2019-2021.} Our target variable are monthly recession indicators for the United States, as defined by the \textit{National Bureau of Economic Research} (NBER). For the empirical replication of the yield curve, we use Treasury yields and macro indicators provided by the \textit{Saint Louis Fed} (FRED).
To ensure a true forecasting horizon, the target is shifted backward by twelve months, such that all predictors available at time $t$ are used to forecast recession status one year ahead at time $t+12$. The yield curve dataset consists of monthly Treasury yields spanning different maturities from three months to thirty years, as released by the US Federal Reserve. These yields are used both directly and as inputs to principal component analysis. 

Macroeconomic indicators include the unemployment rate, year-over-year CPI inflation, year-over-year industrial production growth, and consumer sentiment.\footnote{Consumer Sentiment as measured by the University of Michigan Survey} To distinguish expectations-driven yield curve movements from term premium effects, the ten-year yield is decomposed into an expectations proxy and a term premium component. An adjusted slope measure is constructed by subtracting the term premium from the long rate prior to computing the slope relative to the two-year yield. This allows direct testing of whether recession signals arise from anticipated policy easing or from non-expectational forces such as quantitative easing and safe-asset demand.
